\section{Datos Extras}
\subsection{Articiones prismáticas en D-H}
En este artículo consideramos que entre el eslabón base y el efector final sólo existen articulaciones revolutas, sin embargo el método D-H también cubre el caso cuando existe justamente una articulación prismática entre el eslabón base y el efector final. Tal articulación prismática se considera como un nodo en la cadena cinemática, y la forma en la que se le asigna un marco de referencia es diferente. Puedes investigar más sobre eso. [me lo dijo Gemini y resúmenes de IA de google, y mi propio entendimiento del tema]
\subsection{¿Los marcos de referencia se asignan a los nodos o a los eslabones?}
Realmente los marcos de referencia se asignan a los eslabones y NO a los nodos. Aunque en los dibujos básicamente colocamos los marcos sobre los nodos, el objetivo de estos marcos es ubicar a los eslabones (y no a los nodos). Los eslabones son cuerpos que ocupan espacio en el mundo real, los nodos son sólo el eje de movimiento relativo entre dos cuerpos, son puntos en el espacio, son intangibles. Un marco de referencia sirve para indicarte dónde empieza un eslabón y/o dónde termina otro. Los marcos son para ubicar eslabones. Entonces lo correcto es decir que los marcos se asignan a los eslabones (y no a los nodos). [me lo dijo Gemini]
\subsection{¿La pinza de un brazo robótico es un eslabón?}
Los eslabones son piezas sólidas que se conectan por articulaciones. Si una pinza de un brazo robótico empieza con un eslabón, entonces sí, la pinza es un eslabón porque es una pieza sólida que tiene una articulación que se conecta a otro eslabón. De hecho, esa pinza es el efector final (último eslabón) y además es la herramienta del efector final.

Sólo quiero mencionar que si después del último nodo (la articulación proximal de toda la cadena cinemática) tenemos un pedazo de tubo sólido largo y en el extremo de éste tenemos una piza (pero no hay otra articulación entre el pedazo de tubo y la pinza, sino que están empotrados), entonces el pedazo de tubo y la pinza son un sólo eslabón (el efector final). Aunque el tubo y la pinza son diferentes, dado que no hay una articulación entre ambos, entonces éstos juntos son el último eslabón (la pinza no es un eslabón diferente en sí misma). 
\subsection{Articulación prismática en pinza}
Respecto a la subsección anterior, si la pinza tiene una articulación prismática, ¿se le puede asignar un marco de referencia a la articulación prismática de la pinza (aunque no haya una articulación revoluta entre la pinza y el tubo)?

Usualmente la articulación prismática de una pinza de una brazo robótico no se considera como nodo en D-H, esto porque el objetivo de D-H es ubicar en el espacio al efector final de un brazo robótico, y para ello basta con ubicar el marco del tubo (que es como el chasis o cuerpo de la pinza). El abrir y cerrar de la pinza no cambia la posición de la herramienta (la pinza) en el espacio, sólo cambia su estado de agarre. La articulación prismática de la pinza se considera un mecanismo de actuación independiente. 

Sin embargo, si quisieras ubicar en el espacio una punta de la pinza, entonces sí se necesitaría asignar un marco de referencia a la articulación prismática de la pinza. Y el método de D-H funcionaría bien para ubicar una punta de la pinza (sería como reinterpretar al efector final como una punta de la pinza en lugar de toda la pinza). D-H no funcionaría para ubicar a las dos pinzas aunque se le asignara un marco a la articulación prismática porque entonces tendríamos dos efectores finales, es decir, tendríamos una cadena cinemática abierta no simple (ya no sería una secuencia lineal de eslabones, sino un árbol). Bueno, ahora que lo pienso, podríamos ejecutar D-H dos veces, una para cada dedo de la pinza. 